\subsection{Problem formulation and notation}
\label{sec:formulation}
Let $\Xb,\Xf:\Omega\to\mathbb{R}$ be a baseline and follow-up CT over a voxel
grid $\Omega\subset\mathbb{Z}^3$. Each scan carries a finite lesion set,
$\Lb=\{\ell^b_1,\dots,\ell^b_{m}\}$ and $\Lf=\{\ell^f_1,\dots,\ell^f_{n}\}$,
where each lesion is a connected foreground component with centroid
$\bm{c}\in\Omega$. The goal is a correspondence relation
$\mathcal{M}\subseteq\Lb\times\Lf$ that need \emph{not} be a one-to-one map:
a \emph{merge} is several baseline lesions sharing one follow-up partner,
a \emph{split} is one baseline lesion with several follow-up partners,
a baseline lesion with no partner has \emph{disappeared}, and a follow-up lesion
with no partner is \emph{new}. We write the matching indicator
$M\in\{0,1\}^{m\times n}$ with $M_{ij}=1\Leftrightarrow(\ell^b_i,\ell^f_j)\in\mathcal{M}$,
augmented by per-side ``no-match'' (dustbin) variables
$u^b\in\{0,1\}^{m}$, $u^f\in\{0,1\}^{n}$.

A registration $\bm\varphi$ maps baseline coordinates into follow-up space, so a
baseline centroid $\bm{c}^b_i$ yields a \emph{propagated} prior
$\tilde{\bm{c}}_i=\bm\varphi(\bm{c}^b_i)$. We treat $\tilde{\bm{c}}_i$ as a noisy
observation of the true follow-up location, never as ground truth.

% ------------------------------------------------------------
\subsection{Registration as an uncertain prior}
\label{sec:registration}
Registration produces a deformation field $\bm\varphi$ aligning $\Xb$ to $\Xf$,
used only to propagate known baseline lesion centroids. The propagated point is
modelled as $\tilde{\bm c}_i = \bm c^{f}_i + \bm\delta_i$, where $\bm c^{f}_i$ is
the (unknown) true follow-up centroid and $\bm\delta_i$ is a registration-error
vector whose distribution is \emph{estimated} from longitudinal data
(\cref{sec:calibration}) rather than assumed. Registration is therefore not a
hidden oracle in the pipeline; its uncertainty is pushed into the training
distribution of the segmenter. \todo{specify Locs contribution, and how $\bm\varphi$ is computed.}

% ------------------------------------------------------------
\subsection{Prompt-aware segmentation}
\label{sec:segmentation}
 
\paragraph{One network, three stages.} The segmenter is a single network trained
in three sequential stages: self-supervised pretraining of the encoder, supervised
prompt-aware segmentation on the full single-timepoint corpus, and a supervised
finetune on the longitudinal subset. The network, prompt encoding, and prompt-noise
model defined below are shared across all three. Only the data, the sampling
weights, and the optimiser change from stage to stage (\cref{tab:segstages}).
 
\paragraph{Network.} The backbone $S_\theta$ is a 3D residual-encoder U-Net, which remains SOTA for 3D medical segmentation.
\todo{cite nnU-Net Revisited for the
SOTA claim.} A propagated prior is rendered into a single positive prompt channel
$H^{+}$ concatenated to the CT, so the input is $[\,\Xf;\,H^{+}\,]$.
 
\paragraph{Prompt encoding.} For a point set $P$, each click is rendered as a ball
that falls off linearly with distance and vanishes at radius $r$,
\begin{equation}
H(\bm v)\;=\;s\,\max_{\bm p\in P}\;\phi_r\!\big(\lVert \bm v-\bm p\rVert\big),
\qquad
\phi_r(d)=\Big(1-\tfrac{d}{r}\Big)_{+}.
\label{eq:heatmap}
\end{equation}
The heatmap is recomputed per patch at patch resolution rather than over the whole
volume. Defaults are $r=2$ vox and $s=0.5$.
 
\paragraph{Calibrated prompt noise.} At training time, each in-patch positive
centroid $\bm c$ is perturbed by a propagation offset
\begin{equation}
\bm\delta\sim\mathcal{N}(\bm 0,\,\operatorname{diag}(\bm\sigma^2)),
\qquad
\bm\delta \leftarrow \bm\delta\cdot\min\!\Big(1,\tfrac{\tau}{\lVert\bm\delta\rVert}\Big),
\label{eq:offset}
\end{equation}
i.e.\ an anisotropic Gaussian jitter with a hard radial cap applied to the sampled
offset. We did not pick $\bm\sigma$ and $\tau$ by hand. They come from a
statistical analysis of the inter-timepoint propagation offset of registered points
on the Longitudinal-CT dataset\todo{Longitudinal-CT citation to follow}: we
propagate baseline centroids to follow-up, measure the residual to the matched
follow-up centroid, and take the per-axis spread as $\bm\sigma$ and a radial cap as
$\tau$ (\cref{sec:calibration}). Current values are $\bm\sigma=(2.75,5.19,5.40)$ vox
and $\tau=34$ vox. The jitter lives in the sampler, so it is applied in every
training stage.
 
\paragraph{Training mixture.} Its prompt set is drawn from two rules, designed to
make the model use priors without trusting them:
\begin{enumerate}[nosep,label=(\roman*)]
  \item \textbf{calibrated positive}: in-patch centroids, each jittered by
  \cref{eq:offset};
  \item \textbf{spurious}: with probability $p_{\mathrm{fp}}$, add $n_{\mathrm{fp}}$
  background ``clicks'' (drawn per patch) at distance $\ge d_{\min}$ from any
  foreground voxel, so the segmenter learns not to chase a misplaced prompt.
\end{enumerate}
A patch is centred on a lesion with probability $p_{\mathrm{fg}}$ and otherwise
sampled uniformly, so many patches contain no lesion. A lesion-free patch that
draws a spurious click is our disappeared-lesion case: a prompt is present but there
is nothing to segment. We never place a click at the former lesion site itself, a
design choice we may revisit. The mixture weights differ between the supervised and
finetune stages (\cref{tab:segstages}): the finetune raises both the rate and count
of spurious clicks to suppress hallucinated lesions when the propagated prompt
drifts onto background. Almost every patch therefore carries at least one click,
and at inference every baseline lesion does, so a truly empty prompt is rare in
training and absent at test.
\todo{state final per-stage mixture weights (table).}
 
\paragraph{Objective.} The two supervised stages minimise the standard compound
voxel loss (Dice $+$ cross-entropy, not the connected-component variant) with deep
supervision across decoder scales $s$,
\begin{equation}
\mathcal{L}_{\mathrm{seg}}
=\sum_{s} w_s\big(\mathcal{L}_{\mathrm{Dice}}^{(s)}+\mathcal{L}_{\mathrm{CE}}^{(s)}\big),
\qquad w_s\propto 2^{-s},
\label{eq:segloss}
\end{equation}
with the weights normalised to sum to one.
\todo{cite nnU-Net for the deep-supervised DC$+$CE baseline.}
 
\paragraph{Stage 1: self-supervised pretraining.} The encoder is pretrained by
masked autoencoding on the merged corpus. A random mask is drawn on the encoder's
bottleneck grid and upsampled to voxel resolution (Spark3D-style), hiding a
fraction $\rho=0.75$ of the volume, and the network reconstructs only the masked
regions under an $\ell_2$ loss. Only the encoder is carried into the supervised
stage; the reconstruction decoder is discarded. Optimiser: SGD at $10^{-2}$ with
cosine warm restarts ($T_0=250$), for $250\times1000$ iterations.
\todo{cite MAE and sll revisited for segmentation for DKFZ}
 
\paragraph{Stage 2: supervised segmentation.} Starting from the stage-1 encoder,
the full network is trained on the merged corpus under \cref{eq:segloss}. Sampling uses
$p_{\mathrm{fg}}=0.67$, $p_{\mathrm{fp}}=0.2$, $d_{\min}=50$ vox and a single
spurious click ($n_{\mathrm{fp}}=1$). Optimiser: SGD (Nesterov, momentum $0.99$,
weight decay $3\times10^{-5}$) under the \emph{stretched-tail} polynomial schedule
(transition $k=188$, reference horizon $250$, exponent $0.9$) that lengthens the
low-LR tail relative to standard poly decay, for $500\times1000$ iterations.
\todo{confirm base LR: $10^{-2}$ (config default) vs $10^{-3}$ (training script).}
 
\paragraph{Stage 3: supervised finetuning on Longitudinal-CT.} The full stage-2
network is warm-started and finetuned on the the Longitudinal-CT training
cases. Two changes target the longitudinal regime. First,
sampling is \emph{stratified by lesion type}: each foreground centroid is matched
to its anatomical type from the dataset metadata and given a sampling weight, with
the categories that dominated the failure modes, and that uniform sampling dilutes
behind lung and liver, manually up-weighted (lymph node $\times4$, soft tissue and
skin $\times3$, skeleton $\times4$, all others $\times1$). Second, the prompt
mixture is re-weighted toward spurious clicks ($p_{\mathrm{fg}}=0.55$,
$p_{\mathrm{fp}}=0.5$, $d_{\min}=30$ vox, $n_{\mathrm{fp}}\in[1,3]$). Optimiser:
AdamW at $10^{-5}$ (weight decay $3\times10^{-5}$, gradient-norm clip $1.0$) with a
plain polynomial schedule, for $500\times1000$ iterations.
\todo{cite SPECTRE (Claessens et al., 2026) for low-LR AdamW finetuning of nnU-Net IF I keep using the adamW model}
 
\paragraph{Validation protocol.} \todo{right now I modelled validation to be similar to what the MAST dataset does, where a lesion will always have a prompt, its still up to debate if thius is what we want, but it makes valdaition easier because there a rea lot of unnatotated lesion in the background, and if we have a model that also catches lesion without a corresponding prompt, we will get a lot of Correct false positives}. Validation mirrors deployment: a prompt is always
present, and the lesion either exists (segment it) or has disappeared (output
nothing). We split validation patches by ground-truth foreground into (A)
lesion-with-prompt and (B) prompt-without-lesion, the latter a controlled fraction
$\approx0.3$, each with a guaranteed spurious click. On (A) we report \emph{macro}
Dice, the mean of per-patch foreground Dice, so that small lymph-node, skin and
skeletal lesions count as much as large ones. On (B) we report a false-positive
rate, the mean predicted-foreground fraction (target $\to0$). Finetune checkpoints
are selected on macro Dice. We use this rather than the standard global (micro)
pseudo-Dice, which is dominated by large lesions and was blind to exactly the
small-lesion recall and false-positive suppression we are after.
 
\begin{table}[t]
  \centering
  \caption{Per-stage training configuration of the segmenter. Batch size $4$ and
  the patch size and spacing follow the shared ResEnc-L plan.}
  \label{tab:segstages}
  \small
  \begin{tabular}{@{}p{2.4cm}p{3.3cm}p{3.5cm}p{3.5cm}@{}}
    \toprule
     & SSL pretrain (MAE) & Supervised & Finetune (d013) \\
    \midrule
    Data            & merged (Dataset999) & merged (Dataset999) & Dataset013 subset \\
    Init            & random              & MAE encoder         & full supervised net \\
    Optimiser       & SGD                 & SGD (Nesterov, $m{=}0.99$) & AdamW (clip $1.0$) \\
    LR              & $10^{-2}$           & $10^{-2}$           & $10^{-5}$ \\
    LR schedule     & cosine warm restart ($T_0{=}250$) & stretched-tail poly ($k{=}188$, ref $250$, exp $0.9$) & poly \\
    Weight decay    & $3\times10^{-5}$    & $3\times10^{-5}$    & $3\times10^{-5}$ \\
    Epochs $\times$ iters & $250\times1000$ & $500\times1000$ & $500\times1000$ \\
    Loss            & masked $\ell_2$     & Dice$+$CE, deep sup.  & Dice$+$CE, deep sup. \\
    $p_{\mathrm{fg}}/p_{\mathrm{fp}}/d_{\min}$ & n/a & $0.67/0.2/50$ & $0.55/0.5/30$ \\
    $n_{\mathrm{fp}}$ & n/a               & $1$                 & $1$ to $3$ \\
    Mask ratio $\rho$ & $0.75$            & n/a                 & n/a \\
    Lesion-type strat. & n/a              & n/a                 & LN$\times4$, skin$\times3$, skel$\times4$ \\
    \bottomrule
  \end{tabular}
\end{table}